\section{Literature Review}\label{literature-review}

This section reviews three strands of literature that inform the present
study: (1) investor sentiment and asset pricing, (2) social media as a
financial information channel, and (3) causal identification in
high-frequency panel settings. The paper sits at their intersection and
contributes a formal causal framework using double debiased machine
learning to a domain that has traditionally relied on linear models.

\subsection{Media Sentiment and Asset
Pricing}\label{media-sentiment-and-asset-pricing}

The theoretical foundations for a sentiment channel in asset prices
relies on the concept that media content drives investor behavior as new
information is acquired. As this information spreads through the market,
it influences trader behavior though bid pricing, selling, and buying.

Early text-based evidence comes from Tetlock (2007), who found that
media pessimism in the Wall Street Journal \emph{Abreast of the Market}
column predicts next-day downward pressure on the Dow Jones Industrial
Average, with the effect fully reversing over the following week. This
demonstrates how information available influences trader behavior.
Social media sentiment is just another form of such information.

\subsection{Social Media as a Financial Information
Channel}\label{social-media-as-a-financial-information-channel}

A second branch of literature examines social media specifically as a
source of price-relevant information transmission, distinct from
traditional news outlets. Antweiler and Frank (2004) analyzed
1.5 million messages posted on Yahoo Finance and Raging Bull message
boards and found that message volume predicts next-day trading volume
and volatility, with a more modest link to returns. Their work
establishes the concept that user-generated financial discussion
contains informational content, but leaves open whether the effect is
economically large or causal.

This topic was explicitly covered in a paper that is the most direct
predecessor to this study. Gu and Kurov (2020) used firm-level daily
Twitter sentiment from a commercial data provider to predict
open-to-open stock returns in a Fama-MacBeth, pooled OLS framework over
January 2015 to February 2017. Their key results are as follows: (i) a
significant positive coefficient on one-day lagged Twitter sentiment
(\(+ {0.136}^{***}\)); (ii) a no-reversal test showing only lag~1 is
significant when lags 1--5 are entered jointly; and (iii) a profitable
long-short strategy with an annualized Sharpe ratio of 3.17. These
findings are interpreted as evidence that Twitter sentiment conveys
firm-specific information that is absorbed within one trading day.

The Gu and Kurov (2020) paper also provides an example of the Bloomberg Twitter
sentiment analysis tool used in formal research. % [R1]: Date inconsistency — Introduction and Data say 2025–2026; fix here to match.
This paper replicates
the Gu and Kurov (2020) specifications on an exteded panel of both the S&P~500 and the Russell~3000 from 2016 through 2026 to both test for the robustness of the intial study, and to examine sentiment sensitivity based on firm size as measured in market capitalization.

\subsection{Contribution to
Literature}\label{contribution-to-literature}

The present paper seeks to expand the understanding of how social media
sentiment influences the market. It takes the research question and
variable constructions from the social media--finance literature of
Antweiler and Frank (2004) and Gu and Kurov (2020) and applies more modern causal
identification strategies. This contributes to the broader debate about
whether investor sentiment has real pricing effects. Three features
distinguish it from prior work:

\begin{enumerate}
\def\labelenumi{\arabic{enumi}.}
\item
  \textbf{Formal causal framework.} The DoubleML PLR estimator relaxes
  linearity assumptions that constrain FE and Fama-MacBeth OLS
  estimates, providing a coefficient that is interpretable as a causal
  partial effect under the assumption that the confounder vector of
  confounders \(X_{\text{it}}\) is sufficiently exogenous.
\item
  \textbf{Reverse causality diagnostics.} Rather than assuming the
  direction of causation, the paper explicitly tests for it through a
  placebo regression and lag-structure analysis. The finding that the
  reverse coefficient (\(+ 2.348\)) exceeds the forward estimate
  (\(- 0.924\)) in absolute value disciplines the causal interpretation
  throughout.
\item
  \textbf{Updated sample period.} Prior Twitter-returns studies
  concentrate on 2011--2017, a period of rapid social media adoption.
  The present sample (2024--2025) captures a more mature informational
  environment in which algorithmic trading, bot activity, and
  information saturation may have attenuated the retail sentiment
  channel. The temporal comparison with Gu and Kurov (2020) is informative about the
  stability of the Twitter effect over time.
\end{enumerate}
